% Generated draft: requires mathematical and source review.
\documentclass[11pt]{article}
\usepackage[margin=25mm]{geometry}
\usepackage{fontspec}
\setmainfont{Latin Modern Roman}
\usepackage{amsmath,amssymb,mathtools,unicode-math}
\setmathfont{Latin Modern Math}
\usepackage{enumitem,needspace,xurl,hyperref}
\setlength{\parindent}{0pt}
\setlength{\parskip}{5pt}
\setlength{\emergencystretch}{3em}
\newcommand{\N}{\mathbb{N}}
\newcommand{\Z}{\mathbb{Z}}
\newcommand{\Q}{\mathbb{Q}}
\newcommand{\R}{\mathbb{R}}
\newcommand{\C}{\mathbb{C}}
\newcommand{\F}{\mathbb{F}}
\newcommand{\A}{\mathbb{A}}
\newcommand{\PP}{\mathbb{P}}
\newcommand{\E}{\mathbb{E}}
\newcommand{\eps}{\varepsilon}
\newcommand{\dd}{\,\mathrm{d}}
\begin{document}
\section*{UnsolvedMath ID: 30002902}
\textbf{Source problem code:} catalogue entry
\subsection{Definitions and mathematical statement}
Let \(k\in\Z\) with \(k>1\), and let \(M\) be a compact oriented manifold of dimension \(4k-1\), with fundamental group \(G=\pi_1(M)\). For \(g\in G\), its order \(\operatorname{ord}(g)\) is the least positive integer \(n\) satisfying \(g^n=1_G\), when such an integer exists. Define
\[
\mathcal O_{\mathrm{fin}}(G)
 =\{n\in\Z:n\geq 2\text{ and there exists }g\in G\text{ with }\operatorname{ord}(g)=n\},
\qquad
N_{\mathrm{fin}}(G)=\lvert\mathcal O_{\mathrm{fin}}(G)\rvert.
\]
Thus distinct orders are counted once each, and the identity is excluded.

The supplied record denotes by \(\mathcal S(M)\) a structure group and describes \(\widetilde{\mathcal S}(M)\) as that structure group modulo classes induced by orientation-preserving self-homotopy equivalences of \(M\). Here an orientation-preserving self-homotopy equivalence means a homotopy equivalence \(f:M\to M\) preserving the orientation, with the relevant boundary conditions still to be specified. The record does not define the version of the structure group or the precise quotient construction. If the intended reduced structure group is an abelian group \(A\), its rank can be interpreted as
\[
\operatorname{rank}(A)=\dim_{\Q}(\Q\otimes_{\Z}A).
\]
This interpretation requires confirmation for the intended structure group.

\textbf{Question as supplied.} For every \(k\in\Z\) with \(k>1\) and every compact oriented manifold \(M\) of dimension \(4k-1\), with \(G=\pi_1(M)\), does
\[
\operatorname{rank}\widetilde{\mathcal S}(M)\geq N_{\mathrm{fin}}(G)
\]
hold, under the intended structure-group and quotient conventions?

\textbf{Unresolved qualifications.} The supplied hypotheses do not specify whether \(M\) is closed, which boundary conditions apply if \(\partial M\neq\varnothing\), or which manifold category and structure-group conventions are intended. Connectedness and the convention for \(\pi_1(M)\) are also unstated. No finiteness assumption on \(\mathcal O_{\mathrm{fin}}(G)\) is supplied; interpreting the displayed comparison for infinite cardinalities therefore also requires confirmation.

The original source question recorded locally is: ``Are the elements of the structure group distinguished by this method actually represented by different manifolds?'' The supplied record does not identify the method or the distinguished elements, and it does not specify whether different manifolds means distinct homeomorphism classes, diffeomorphism classes, or another equivalence relation. This original question is retained as a related target; its equivalence to the displayed rank inequality is not established by the supplied evidence.
\subsection{Short English statement}
The supplied question asks whether the rank of a reduced structure group of every compact oriented \((4k-1)\)-manifold, for \(k>1\), is at least the number of distinct nontrivial finite element orders in its fundamental group. The structure-group conventions, boundary conditions, and relation to the original question about distinct manifolds require clarification.
\subsection{Sources}
\begin{itemize}
\item Supplied local UnsolvedMath registry record, permanent ID \(30002902\), title \emph{Lower Bounds for Reduced Structure Groups}, exact source problem code \texttt{catalogue entry}; collection: Oberwolfach Reports Open Problems. No individual proposer is identified. The supplied proposed year is 2015.
\item Supplied source citation: \emph{Noncommutative Geometry (2016)}, Oberwolfach Reports, DOI \texttt{10.4171/owr/2015/29}, pp.~1631--1700. \url{https://doi.org/10.4171/owr/2015/29}. The supplied report-title year is 2016, while the DOI and proposed-year metadata contain 2015; this discrepancy is retained without resolution.
\item Supplied literature reference: Weinberger--Yu, \emph{Geometry and Topology}. \url{https://msp.org/gt/2015/19-5/gt-v19-n5-p08-p.pdf}. The record provides no article title or precise theorem reference.
\item Supplied literature reference: Oberwolfach report 29. \url{https://doi.org/10.4171/owr/2015/29}.
\item Literature assessment dated exactly \texttt{2026-08-21}: the supplied curation labels the problem open and reports that the rank lower bound is proved under additional residual-finiteness and torsion hypotheses, while its general validity and the claim that the detected elements represent distinct manifolds remain conjectural. The precise additional hypotheses are not supplied. This is supplied dated curation, not an independently verified literature assessment.
\item The record labels its statement \texttt{corrected\_verified} and reports that it was checked against the official Oberwolfach report and preserves that report's hypotheses, notation, and target. These are supplied source-verification claims; no independent source consultation or verification is asserted here.
\end{itemize}
\end{document}
